\section{Analysis and Falsification}

\subsection{Event-phase analysis}

Repeated EEG windows are nested within events, sessions, and participants. Mixed-effects or other preregistered repeated-measure models estimate within-child spectral change across preswitch reference, challenge, regulation, recovery monitoring, and stable return.

Current rule, switch direction, task difficulty, sensory load, strategy type, motion burden, session number, time on task, acquisition profile, and signal quality are prespecified contextual variables.

\subsection{Behavioral-neural recovery coupling}

The core behavioral endpoint \(T_{\mathrm{rec}}\) remains independently defined from task behavior. \(T_{\mathrm{EEG}}\) is constructed without using behavioral recovery or external clinical outcomes to tune its threshold.

The analysis estimates association and uncertainty between \(T_{\mathrm{EEG}}\) and \(T_{\mathrm{rec}}\), allowing repeated events and censoring. Concordance is evidence of mechanistic coherence, not proof that one process causally generates the other.

\subsection{Incremental information}

The EEG channel is retained scientifically only if it improves a prespecified participant-held-out prediction or decision metric beyond behavioral and contextual features alone. The analysis must compare identical splits and outcome definitions.

Important falsification patterns include no reproducible phase effect, unstable neural recovery estimates, disappearance of apparent effects after motion/EMG control, weak within-child repeatability, no behavioral-neural coupling, and no incremental out-of-sample value.
